% TOP_PROBLEM: {"id": 4800031, "title": "Multiple ergodic averages — Problem 31", "category_id": 11, "category": {"id": 11, "name": "dynamical_systems", "display_name": "Dynamical Systems", "description": "Problems about long-term behavior of deterministic systems, Hamiltonian dynamics, and periodic orbits.", "slug": "dynamical-systems", "order_index": 11, "created_at": "2026-07-31T15:26:25.671Z"}, "status": "open", "problem_number": "AMR-047-0031", "set_id": 13}
% FIRST_POSTED: 2026-10-01
% ENTRY_KIND: statement_only
% UnsolvedMath ID 4800031; source catalogue code AMR-047-0031
% !TeX program = lualatex
% Definition and sources only; this file is not an LLM solution attempt.
\documentclass[11pt,a4paper]{article}
\usepackage[margin=25mm]{geometry}
\usepackage{fontspec}
\setmainfont{lmroman10-regular.otf}[BoldFont=lmroman10-bold.otf,
 ItalicFont=lmroman10-italic.otf,BoldItalicFont=lmroman10-bolditalic.otf]
\usepackage{amsmath,amssymb,mathtools}
\usepackage{unicode-math}
\setmathfont{latinmodern-math.otf}
\AtBeginDocument{\let\setminus\smallsetminus}
\usepackage{xurl,hyperref}
\hypersetup{colorlinks=true,urlcolor=blue}
\setcounter{secnumdepth}{0}
\setlength{\parindent}{0pt}
\setlength{\parskip}{5pt}
\setlength{\emergencystretch}{3em}
\begin{document}
\section*{Multiple ergodic averages — Problem 31}\label{4800031}
{\small UnsolvedMath ID 4800031 · AMR-047-0031 · Dynamical Systems · AMR Open Problem Lists}
\subsection{Definitions and mathematical statement}
Let $(\sigma_k)_{k\ge1}$ be a nonincreasing
sequence in $(0,1]$ with $k\sigma_k\to\infty$.
On a probability space $(\Omega,\mathcal F,\mathbb P)$,
take independent Bernoulli variables $\xi_k$
with $\mathbb P(\xi_k=1)=\sigma_k$.
Almost surely the selected set is infinite;
let $a_n(\omega)$ be its $n$th smallest member.

Prove that there is one event $\Omega_0$ of
probability one such that for every
$\omega\in\Omega_0$, every $\ell\ge1$, every
probability-preserving commuting invertible system
$(X,\mathcal B,\mu,T_1,\ldots,T_\ell)$, and all
$f_i\in L^\infty(\mu)$,
\[
 \frac1N\sum_{n=1}^N\prod_{i=1}^\ell f_i\circ T_i^{a_n(\omega)}
 \quad\text{and}\quad
 \frac1N\sum_{n=1}^N\prod_{i=1}^\ell f_i\circ T_i^n
\]
have the same $L^2(\mu)$ limit. Also prove
that each such selected sequence gives multiple
recurrence: every $A$ with $\mu(A)>0$ admits
some $n$ with
$\mu(A\cap T_1^{a_n(\omega)}A\cap\cdots
\cap T_\ell^{a_n(\omega)}A)>0$.

The full-probability event must be chosen before
the dynamical system and observables are chosen;
an event depending on each tested system would
be weaker. Averages are normalized by the
number of selected times $N$, not by the
largest selected integer. No ergodicity or
pointwise convergence is imposed. The sparsity
condition and monotonicity concern the selection
probabilities, not deterministic gaps between
successive selected integers.
\subsection{Short English statement}
Does sufficiently frequent independent random sampling almost surely preserve every commuting multiple ergodic mean and recurrence property?
\subsection{Sources}
\begin{itemize}
\item[S1] ULAM.AI and the UnsolvedMath Contributors, supplied problems.json, record 4800031 (AMR-047-0031), AMR Open Problem Lists. Catalogue identity and source transcription; CC BY 4.0.
\url{https://huggingface.co/datasets/ulamai/UnsolvedMath}.
\item[S2] Frantzikinakis - Some open problems on multiple ergodic averages (2016), Problem 31. Original source indexed by AMR-047-0031.
\url{https://arxiv.org/abs/1103.3808}.
\end{itemize}
\end{document}
